\documentclass[10pt]{article}

\usepackage[utf8]{inputenc}

\usepackage[T1]{fontenc}

\usepackage{amsmath}

\usepackage{amsfonts}

\usepackage{amssymb}

\usepackage[version=4]{mhchem}

\usepackage{stmaryrd}

\usepackage{bbold}



\begin{document}

Для обозначения Р. ч. $r$ применяются рациональные дроби $\frac{a}{b}$ из класса эквивалентности, задающего это число: $\frac{a}{b} \in r$. Таким образом, одно и то же Р. ч. может быть записано разными, но эквивалентными рациональными дробями



Если каждому Р.ч., содержащему рациональную дробь вида $\frac{a}{1}$, поставить в соответствие целое число $a$, то получится изоморфное отображение множества указанных Р. ч. на кольцо $\mathbb{Z}$ целых чисел. Поэтому Р. ч., содержащие рациональные дроби вида $\frac{a}{1}$, обозначают через $a$.



Всякая функция вида



\[

\varphi(r)=|r|^{\alpha}, \quad 0<\alpha \leqslant 1,

\]



является метрикой в поле Р.ч. $Q$, то есть удовлетворяет условиям:



\begin{enumerate}

  \item $\varphi(r)>0$ при любом $r \neq 0, \varphi(0)=0$;



  \item $\varphi\left(r^{\prime}+r^{\prime \prime}\right) \leqslant \varphi\left(r^{\prime}\right)+\varphi\left(r^{\prime \prime}\right)$;



  \item $\varphi\left(r^{\prime}, r^{\prime \prime}\right)=\varphi\left(r^{\prime}\right) \varphi\left(r^{\prime \prime}\right)$



\end{enumerate}



при любых $r^{\prime} \in \mathbb{Q}$ и $r^{\prime \prime} \in \mathbb{Q}$. Поле Р. ч. не является полным в метрике (1). Пополнением поля Р.ч. по метрике (1) является поле действительных чисел.



Функция



\[

\Psi_{p}(r)=\rho^{v(r)}

\]



где $p$ - простое число, $r$-Р.ч., представшмое в виде



\[

r=p^{v(r)} \frac{a}{b}

\]



$\left(v(r)\right.$ - целое, $\frac{a}{b}$ - несократимая рациональная дробь, причем числа $a$ и $b$ не делятся на $p$ ), а $\rho$ - фиксированное число, $0<\rho<1$, также является метрикой в поле Р. ч. $\mathbb{Q}$. Она наз. $p$-а д и ч е с к о й м е т р ик о й. Поле Р. ч. $\mathbb{Q}$ с метрикой (2) не является полным. Пополнением поля Р.ч. по метрике (2) является поле p-адических чисел. Метрики (1) и (2) (для всех простых чисел) исчерпывают все нетривиальные метрики в поле P. ч.



В десятичной записи Р.ч. и только они гредставимы периодическими десятичными дробями.



Лum.: [1] Б о ре в ич З. И., ІІ а фа р е в и ч И. Р., Теория чисел, 2 изд., М., 1972; [2] Пн и з о Ш., З а м а н с к и й М.,



РАЩИОНАЛЬНОСТИ ТЕОРЕМЫ Д. Д. Кудряөчөв. р а и ч е с к их г р у п п- утверждения о рациональности (унирациональности) или нерациональности тех или иных групповых алгебраич. многообразий. Так как абелевы многообразия всегда нерациональны, то основной интерес представляют Р. т. для линейных алгебраич. групп, Здесь проблема рациональности имеет два существенно различных аспекта: геометрическиї и арифметический, отвечающие соответственно алгебраически замкнутому и незамкнутому основному полю $K$. Первые Р. т. над полем $\mathbb{C}$ комплексных чисел были фактически доказаны еще Э. Пикаром (E. Picard) и в современной терминологии устанавливают унирациональность многообразий связных комплексных групі. В явной форме проблема рациональности групповых многообразий была поставлена К. IIевалле [1] лишь в 1954. Прогресс в этом направлении тесно связан с достижениями структурной теории алгебраич, групп. Так, разложение Леви позволяет свести проблему рациональности к редуктивным группам, а разложение Брюа - доказать рациональность многообразий редуктивных групп над любым алгебраически замкнутым полем. Таким образом, в геометрич. случае имеется окончательный результат. Гораздо более сложной оказывается ситуация для незамкнутых полей $K$. Примеры нерациональных $K-$ многообразий доставляют уже алгебраич. торы; напр., трехмерный норменный тор $T=R_{L / K}^{(1)}\left(G_{m}\right)$, соответствующий биквадратичному расширению $L=K(\sqrt{a}, \sqrt{b})$ поля $K$ (см. [1]). Этот пример минимален, ибо торы размерности $\leqslant 2$ рациональны. В общем случае алгебраич. торы всегда унирациональны. Произвольные связные $K$-группы не обязательно унирациональны [3], однако, если поле $K$ совершенно или группа $G$ редуктивна, унирациональность имеет место (см. [1] - [4]). Тем самым проблема рациональности групповых многообразий имеет характер Люрота проблемы над незамкнутым полем.



Так как произвольная редуктивная группа является почти прямым произведением тора и полупростой группы, то естественно различать два основных случая: 1) $G$ - тор; 2) $G$ - полупростая группа. В первом случае исследование проводится при помощи различных когомологич. инвариантов (для полупростых групп әти инварианты оказываются неәффективными). Достаточно законченные результаты имеются для торов, разложимых над абелевым расширением поля определения (см. [5]). Первый пример нерационального многообразия в классе полупростых групп был неодносвязной группой, конструкция к-рой фактическІ содержится в [10]. Возникшая при этом гипотеза о том, что многообразия односвязных групп всегда рациональны, была решена отрицательно В. П. Платоновым при помощи развитой им п р и в е д е н н о й $K$ - т е о р и и (см. [6], [7]). Оказалось, что приведенная групиа Уайтхеда $S K_{1}(D)$ конечномерной центральной простой $K$-алгебры $D$ тривиальна, если многообразие, определяемое $S L(1, D)$, рационально над $K$. Эти результаты были перенесены на унитарные группы [12]. Ряд результатов связан с исследованием рациональности спинорных многообразий $\operatorname{Spin}(n, f)$, где $f$ - невырожденная квадратичная форма над $K$ от $n$ переменных (char $K \neq 2$ ). Спинорные многообразия рациональны, если либо $n \leqslant 5$, либо поле $K$ является недискретным локально компактным или полем рациональных чисел (см. [8], [9], [11]); для $n \geqslant 6$ суцествуют нерациональные спинорные многообразия [8]. Последний результат удивителен тем, что Spin $(n, f)$ является двулистным накрытием рационального многообразия $\mathrm{SO}(n, f)$.



Термин «Р. т.» иногда употребляется в теории алгебраич. групп в несколько ином смысле, применительно к утверждениям о свойствах групп над не обязательно алгебраически замкнутым полем. $\mathrm{K}$ утверждениям такого типа относится, напр., т е о р е м а $\mathrm{P}$ оз е нли х т а - Г р о те нд и к а о том, что любая связная $K$-группа обладает максимальным тором, определенным над $K$ (см. [4]).



Jum.: [1] C h e v a 11 e y C., "J. Math. Soc. Japan", 1954, v. 6, № 3/4, p 303-24; [2] D e m a z u r e M., G r o t h e n d il i c h t M., "Ann. mat. pura appl.», 1957, v. 43, p. 25-50; [4] Б о р е л ь. А., Јинейные алгебраические группы, пер. с англ. М., $1972 ;$ [5] , В о с к р е с е н с к и й В. Е., Алгебраические торы, М., 1977; [6] П Л а т о н о в В. П., "Тр. Матем. ин-та АН 1977, т. 21 , № 3, c. 197-98; [8] е г o ж е, "Докл. АH CCCP", 1979, т. 248, № 3 , с $524-27 ;$ [9] е г о ж е, "Тр. Матем. ин-та AH CCCP》, 1981 , т. 157, c. $161-69$; [10] C e p p Ж - П., Когомологии Галуа, пер. с франц., М., 1968; 111 с к и й В. И., «Матем. со́.», 1979, т. 110, № 4, с. 579-96.



РЕАЛИЗУЕМОСТЬ - один из видов неклассич. интерпретаиий логических и логико математических языков. Различные интерпретации типа Р. определяются по следующей схеме. Для формул логико-математич. языка определяется отношение «объект $e$ реализует замкнутую формулу $F_{»}$, к-рое сокращенно записывается «erF». Определение носит индуктивный характер:





\end{document}